\documentclass[11pt]{article}

% ============================================================
%  Direct collider bounds on KK masses in warped (RS) models
%  A slow, pedagogical internal review.
%  Tracks (a) the collider physics, (b) what OUR code does,
%  (c) the subtleties and the gap to "rigorous".
%  Internal working note, not for distribution.
% ============================================================

\usepackage[margin=1in]{geometry}
\usepackage{amsmath,amssymb,amsthm}
\usepackage{mathtools}
\usepackage{graphicx}
\usepackage{booktabs}
\usepackage{enumitem}
\usepackage{xcolor}
\usepackage{framed}
\usepackage[colorlinks=true,linkcolor=blue,citecolor=blue,urlcolor=blue]{hyperref}

% lightweight "callout" box
\colorlet{shadecolor}{gray!12}
\newenvironment{callout}[1]{%
  \par\medskip\begin{shaded}\noindent\textbf{#1}\par\smallskip}{\end{shaded}\medskip}

\newcommand{\MKK}{M_{\mathrm{KK}}}
\newcommand{\LIR}{\Lambda_{\mathrm{IR}}}
\newcommand{\sw}{s_W}
\newcommand{\cw}{c_W}
\newcommand{\swsq}{s_W^2}
\newcommand{\cwsq}{c_W^2}
\newcommand{\Pibig}{\Pi}
\newcommand{\half}{\tfrac{1}{2}}
\newcommand{\vev}{v}
\newcommand{\eps}{\varepsilon}
\newcommand{\SUL}{SU(2)_L}
\newcommand{\SUR}{SU(2)_R}
\newcommand{\UX}{U(1)_X}
\newcommand{\PLR}{P_{LR}}
\newcommand{\Tparam}{\hat{T}}
\newcommand{\gKK}{g_{\mathrm{KK}}}
\newcommand{\GKK}{G_{\mathrm{KK}}}
\newcommand{\sBR}{\sigma\!\times\!\mathrm{BR}}
\newcommand{\code}[1]{\texttt{#1}}
\newcommand{\todo}[1]{{\color{red}[#1]}}

\theoremstyle{definition}
\newtheorem{definition}{Definition}

\title{\bf Direct Collider Bounds on Kaluza--Klein Masses in Warped Models\\[4pt]
\large A slow, pedagogical internal review --- physics, code, and subtleties}
\author{Internal working note (5D--Neutrino--Mixing repo)}
\date{\today}

\begin{document}
\maketitle

\begin{abstract}
\noindent
This note is a from-scratch walk-through of the \emph{direct} collider searches
that bound the lightest Kaluza--Klein (KK) mass in a Randall--Sundrum (RS)
warped extra dimension, and of how our catalogue encodes them as the
\code{collider\_rs} family (\code{CR001}--\code{CR014}). It assumes no prior
knowledge of resonance searches, vector-like quarks, or KK gravitons --- each is
built up slowly. Part~\ref{part:why} explains why a fourteen-process direct-search
family belongs in a flavour/EW catalogue at all, and why these particular bounds
are the \emph{least} binding ones we carry. Part~\ref{part:phys} defines the
collider observables: resonant production of a heavy state, its decay channels,
the cross-section-times-branching-ratio ($\sBR$) curve, and the mass-exclusion
edge a search reports. Part~\ref{part:rs} explains how RS produces each kind of
signal --- KK gluon, KK graviton, KK electroweak vectors, and the custodial
vector-like quark (VLQ) partners --- and where a near-flat $\MKK$ floor comes
from. Part~\ref{part:model} addresses custodial and alignment dependence and
states plainly that \emph{our treatment is a mass-exclusion proxy, not a
rigorous recast}. Part~\ref{part:code} documents exactly what the code does
(the shared adapter \code{collider\_resonance.py}, the per-process
\code{CR0xx.py}, their YAML sidecars, severities, and the quark-only allowlist
split). Part~\ref{part:subtle} catalogues the subtleties and the precise gap to
``rigorous,'' including the $\MKK$ convention. An annotated reading guide closes
the note.
\end{abstract}

\tableofcontents

% ====================================================================
\part*{Orientation: why this note exists}
\addcontentsline{toc}{section}{Orientation: why this note exists}
\label{part:why}
% ====================================================================

Our constraint catalogue is dominated, numerically, by \emph{indirect} probes:
meson mixing ($\eps_K$, $\Delta m_s$, $\phi_s$), the $Z\to b\bar b$ vertex, and
the oblique parameters. Those are low-energy observables that feel a heavy KK
tower only through tiny virtual effects, and precisely because the effects are
tiny, even a small measured deviation translates into an enormous lower bound on
$\MKK$ --- tens of TeV in the anarchic-flavour regime. The \code{collider\_rs}
family is the \emph{opposite} kind of constraint: a heavy KK state is produced
\emph{on shell} at the LHC and seen (or excluded) directly. These direct
searches are conceptually the cleanest bounds we have --- ``did we make the
particle or not'' --- yet in our scan they are the \emph{weakest}, because the
energy of even a 13~TeV collider caps the reach at a few TeV, whereas the
indirect bounds already sit far above that. The reason for the cap is the
\emph{parton luminosity}: to make a resonance of mass $M$ the colliding partons
must carry a fraction $x\sim M/\sqrt{s}$ of each proton's momentum, and the
parton densities fall steeply with $x$. At $\sqrt s = 13$~TeV a 6~TeV resonance
already needs $x\sim 0.5$, where the quark/gluon densities are tiny, so the
production cross section collapses long before $\sqrt s$ is reached. That is why
the cleanest bounds are also the weakest. Understanding this, and exactly how
honestly we encode it, is the purpose of this note.

A one-paragraph executive summary, to be unpacked over the following pages:

\begin{quote}
\small
A warped extra dimension predicts a whole spectrum of TeV-scale resonances: a
colour-octet KK gluon that decays overwhelmingly to top pairs, a spin-2 KK
graviton that shows up in diboson and diphoton channels, KK partners of the
electroweak gauge bosons, and (in the custodial version) vector-like quark
partners of the top and bottom. The LHC hunts every one of these as a bump in an
invariant-mass spectrum, or as pair-produced heavy quarks, and reports a
mass-exclusion edge: ``no signal below $M_\star$.'' Collected together, these
edges form a near-flat \emph{direct-search floor} on $\MKK$ of order a few TeV.
Our catalogue carries fourteen of them (\code{CR001}--\code{CR014}), but it does
\emph{not} run a real collider recast: each one is a documented
\emph{mass-exclusion proxy} that interprets the scan's single KK mass scale as
the relevant resonance mass and asks only whether it clears the catalogued
benchmark edge. The proxy is honest about being a proxy --- every module is
tagged \code{NEEDS-HUMAN-PHYSICS} --- and because the rigorous flavour bounds
already push $\MKK$ above $\sim 25$~TeV (physical), these few-TeV direct floors
essentially never bite once the indirect bounds are switched on.
\end{quote}

\noindent The logical spine of the whole note, in one picture:
\begin{footnotesize}
\begin{verbatim}
  RS spectrum at the TeV scale
       |
       +-- KK gluon g_KK^(1)  --(decays to)-->  t tbar        ==> CR001  (resonance bump)
       |
       +-- KK graviton G_KK^(1) (spin 2) ------>  WW/ZZ, gamma gamma  ==> CR007, CR013
       |
       +-- KK EW vectors (spin 1)  ----------->  WW/WZ/ZZ, ll, l nu   ==> CR012; CR005,CR006(OUT)
       |
       +-- custodial VLQ partners (T_5/3, T, B, (T,B))  -- pair produced --> CR002/3/4/8/10
       |
       +-- contact / VBS / four-top tails  ------------------------------> CR009,CR011,CR014(OUT)
                           |
   each search => an experimental MASS-EXCLUSION EDGE M*  (a few TeV)
                           |
   OUR CODE: m_predicted = (a KK mass scale)  ;  pass iff  m_predicted >= M*
             ratio = M* / m_predicted  ;  ratio <= 1 passes   [PROXY, not a recast]
                           |
   net effect: near-flat direct floor ~ few TeV
               -> SMOTHERED by the rigorous flavour floor (~25-30 TeV physical)
               -> at M_KK = 1 TeV CR001 excludes; by ~30 TeV all CR* clear
\end{verbatim}
\end{footnotesize}
\noindent\emph{Legend (all defined where each process is introduced in
Part~\ref{part:rs}):} $T_{5/3},T,B,(T,B)$ are the custodial vector-like quark
partners (\S\ref{sec:vlq}); ``HVT model~B'' is the heavy-vector-triplet diboson
benchmark for the KK spin-1 vector (\S\ref{sec:ewkk}); ``VBS'' is vector-boson
scattering and ``INFO'' marks the lone non-vetoing record (\code{CR011},
\S\ref{sec:severity}).

\noindent Keep this map in mind; every branch is unpacked below. The single most
important honesty point --- defused here at the top and revisited throughout ---
is that the arrow ``$\sBR$, width, interference, acceptance'' that a real recast
would compute is \emph{not} in our code. We compare masses, not rates.

% ====================================================================
\part{The collider physics, from scratch}
\label{part:phys}
% ====================================================================

\section{What a ``direct search'' actually measures}
\label{sec:directsearch}

Imagine some new heavy particle $X$ of mass $M_X$ a few TeV. At a hadron
collider two protons collide; their constituent quarks and gluons can, with some
probability, fuse to make $X$. If $X$ is unstable it decays promptly into
ordinary particles. A ``direct search'' is the experimental program of looking
for the kinematic signature of that production-and-decay chain in the recorded
events. There are two archetypes, and every \code{CR0xx} constraint is one of
them.

\emph{(Where do the heavy states $X$ come from? A warped extra dimension predicts
a whole tower of TeV-scale KK resonances --- a colour-octet KK gluon, a spin-2 KK
graviton, KK partners of the electroweak gauge bosons, and custodial vector-like
quark partners of the top and bottom. The full RS spectrum is built up in
Part~\ref{part:rs}; for now it is enough to know that such states exist and are a
few TeV heavy. Terms like ``vector-like quark,'' ``top-philic,'' and ``boosted
top'' that appear below are defined where each process is introduced in
Part~\ref{part:rs}.)}

\begin{definition}[Resonance (bump) search]
A single heavy state $X$ is produced and decays to a definite final state
$f\bar f$ or $VV$. Because $X$ has a definite mass, the invariant mass of its
decay products $m_{f\bar f}$ piles up in a narrow ``bump'' centred at $M_X$,
sitting on a smooth Standard-Model (SM) background. The search scans $m_{f\bar
f}$, sees no bump, and reports an upper limit on the production rate
$\sigma(pp\to X)\times \mathrm{BR}(X\to f\bar f)$ as a function of $M_X$.
\end{definition}

\begin{definition}[Pair-production search]
A heavy fermion $Q$ (here a vector-like quark) is produced in pairs $Q\bar Q$
through its (model-independent) strong coupling to gluons. Each $Q$ then decays
to a SM quark plus a $W$, $Z$, or Higgs. The search demands the full cascade
($b W$, $t Z$, $t H$, $\dots$) and reports a lower limit on $M_Q$ for an assumed
mixture of branching fractions. Because the production is QCD-driven, the bound
is nearly model-independent in the production cross section and depends mainly on
the assumed decay pattern.
\end{definition}

In both cases the quantity the experiment actually hands you is a curve, not a
number: an upper limit on $\sBR$ versus mass (resonance), or equivalently a
\emph{mass-exclusion edge} $M_\star$ where the model's own predicted $\sBR$ curve
crosses the experimental limit. Below $M_\star$ the model is excluded; above it
the rate is too small to see. The edge is therefore \emph{benchmark-dependent}:
it presumes a specific signal model (production coupling, width, branching
fractions). This benchmark dependence is the crack through which all of our
proxy approximations enter (\S\ref{part:model}).

\section{Cross section, branching ratio, width: the three ingredients}
\label{sec:sbr}

The rate of a resonance signal factorises, schematically, as
\begin{equation}
N_{\mathrm{signal}} \;\propto\;
\underbrace{\sigma(pp\to X)}_{\text{production}}\;\times\;
\underbrace{\mathrm{BR}(X\to f\bar f)}_{\text{decay fraction}}\;\times\;
\underbrace{A\,\epsilon}_{\text{acceptance}\times\text{efficiency}},
\label{eq:rate}
\end{equation}
and the comparison to data is done as a function of $M_X$. Three model inputs
control it:
\begin{itemize}[leftmargin=1.4em]
\item \textbf{Production cross section} $\sigma(pp\to X)$: set by how strongly $X$
couples to the colliding partons (light quarks and gluons) and by the parton
densities, which fall steeply with mass --- this is the ``parton luminosity''
suppression that caps the reach at high mass.
\item \textbf{Branching ratio} $\mathrm{BR}(X\to f\bar f)$: the fraction of
decays into the searched channel. An RS KK gluon, for instance, is
\emph{top-philic} --- it couples most strongly to the IR-localised top (the RS
localisation mechanism is recalled in Part~\ref{part:rs}), so
$\mathrm{BR}(\gKK\to t\bar t)$ is large and the $t\bar t$ channel is the place
to look.
\item \textbf{Total width} $\Gamma_X$: how broad the bump is. A narrow resonance
gives a sharp peak; a broad one ($\Gamma_X/M_X\sim 0.1$--$0.5$) smears across
many mass bins and is much harder to distinguish from background. The RS KK gluon
is famously \emph{broad}: the canonical Lillie--Randall--Wang value is
$\Gamma/M\sim 0.17$ (configuration-dependent, $\sim 0.15$--$0.3$ across the
literature), and our sidecar records the benchmark used here as
\code{width\_over\_mass: 0.153}.
\end{itemize}
Interference between the resonance and the SM continuum, and the detector
acceptance, complete the list. \textbf{A rigorous recast computes all of these
from the model point.} Our code computes \emph{none} of them; it uses only the
mass. That is the entire content of the ``proxy'' tag, and it is worth stating
this early and bluntly.

\section{The effective-operator view (for the tails)}
\label{sec:eft}

Two of the fourteen searches are not bump hunts at all but \emph{tail}
measurements, and they are cleanest in effective-field-theory (EFT) language.
When the resonance is too heavy to produce on shell, integrating it out leaves a
\emph{contact operator} --- a four-fermion interaction suppressed by a scale
$\Lambda$. The classic example is high-mass Drell--Yan: a virtual KK photon/$Z$
or KK gluon exchanged between quarks and leptons generates
\begin{equation}
\mathcal{L}_{\mathrm{contact}} \;\sim\; \frac{4\pi}{\Lambda^2}\,
(\bar q\,\Gamma\, q)(\bar\ell\,\Gamma\,\ell),
\qquad
\Lambda \;\sim\; \frac{\MKK}{\sqrt{|g_q g_\ell|}},
\label{eq:contact}
\end{equation}
which distorts the smooth high-mass tail of the dilepton spectrum rather than
producing a peak. The search bounds $\Lambda$; our \code{CR009} pins, by an
explicit \emph{conservative} policy (\code{CR009.py} selects the
\emph{weakest} catalogued edge), the destructive-interference limit
$\Lambda_{LL} > 23.9$~TeV from CMS. Stronger benchmarks (e.g.\ the ATLAS
constructive $\Lambda_{LL}^{+} > 35.8$~TeV) are retained only as diagnostics. The longitudinal vector-boson
scattering (VBS) measurement (\code{CR011}) is the analogous probe of the
\emph{strong electroweak-symmetry-breaking} sector: it tests whether $W_L W_L$
scattering grows with energy as it would if the Higgs sector were composite, and
reports a fiducial cross-section upper limit. These tail observables bound the
\emph{scale} $\Lambda$ or a fiducial rate, not a resonance mass --- a distinction
that matters for how the proxy treats them (\S\ref{part:code}).

% ====================================================================
\part{How RS generates each signal, and where the floor comes from}
\label{part:rs}
% ====================================================================

\section{RS recap: the spectrum at the TeV scale}
\label{sec:rsrecap}

Recall the geometry (repo conventions). Spacetime is a slice of AdS$_5$ with a
UV brane at the Planck scale and an IR brane at the TeV scale; the warp factor
$\eps = e^{-k\pi r_c}\sim 10^{-15}$ generates the hierarchy. Every bulk field
has a tower of KK modes whose masses are set by the IR scale,
\begin{equation}
m_n \;=\; x_n\,\LIR, \qquad
x_1^{\rm gauge} \approx 2.45, \quad
x_1^{\rm grav} \approx 3.83,
\label{eq:KKmass}
\end{equation}
where the $x_n$ are roots of a Bessel function and $\LIR$ is the IR scale. The
root depends on the spin: the first KK \emph{gauge} mode sits near the first zero
of $J_0$ ($2.405$ in the $\eps\to0$ Neumann limit, shifted to the finite-$\eps$
NN root $x_1^{\rm gauge}\approx 2.45$), whereas the first KK \emph{graviton}
mode sits at the first zero of $J_1$, $x_1^{\rm grav}\approx 3.83$, so the
lightest graviton is heavier than the lightest gauge resonance by a factor
$x_1^{\rm grav}/x_1^{\rm gauge}\approx 1.5$ (Agashe--Davoudiasl--Perez--Soni,
arXiv:hep-ph/0701186). The physical first KK mass is therefore
$\MKK^{\mathrm{phys}} = x_1\LIR$, with $x_1$ the spin-appropriate root --- a
factor we will have to keep track of carefully (\S\ref{sec:Mconvention}), because
collider papers quote masses of \emph{physical} resonances while some theory
formulas use $\MKK\equiv\LIR$.

What sits in the tower? The bulk gauge fields give KK gluons (colour octet, the
heaviest-produced and most strongly coupled), KK $W$/$Z$/$\gamma$ (electroweak
vectors), and --- because RS is a theory of \emph{gravity} in 5D --- KK
gravitons (spin 2). Bulk fermions give KK quark and lepton partners; in the
custodial model (\S\ref{part:model}) the third-generation partners are
\emph{vector-like quarks} with exotic and standard charges. Each of these is a
collider target, and our fourteen processes enumerate them.

\section{The KK gluon \texorpdfstring{$\to t\bar t$}{-> ttbar} (CR001)}
\label{sec:cr001}

The KK gluon is the flagship RS collider signature. It is a colour octet, so it
is produced in $q\bar q$ annihilation through the QCD (colour) interaction --- but
crucially \emph{not} at full strong-coupling strength. The light quarks that
dominate the proton's parton luminosity are UV-localised, and their overlap with
the IR-peaked KK-gluon wavefunction is small, so the production coupling to light
quarks is \emph{suppressed} to roughly $0.2\,g_s$ (Lillie--Randall--Wang,
arXiv:hep-ph/0701166). Gluon fusion $gg\to\gKK$ is moreover \emph{forbidden} by
wavefunction orthogonality (the flat SM-gluon zero mode is orthogonal to the
first KK mode), so the dominant gluon parton luminosity does not contribute at
all. This suppressed light-quark coupling is the main reason the KK-gluon
production cross section falls and the direct reach caps at a few TeV. By
contrast the KK gluon couples \emph{strongly} to whatever is localised near the
IR brane: in anarchic RS the top quark is the most IR-localised SM fermion, so
the KK gluon is \emph{top-philic}, $\mathrm{BR}(\gKK\to t\bar t)$ dominates, and
the relevant search is a $t\bar t$ resonance hunt. (Lillie--Randall--Wang made
exactly this point and motivated boosted-top reconstruction; it is the paper-era
theory reference in our \code{CR001.yaml}.)

Experimentally one looks for a broad bump in the $t\bar t$ invariant-mass
spectrum. The current best benchmark exclusion we encode is the full-Run-2 CMS
all-topology search \mbox{CMS-B2G-25-009} (arXiv:2603.23454), which excludes the
RS KK-gluon benchmark over $0.5 < m(\gKK) < 5.5$~TeV; our \code{CR001.yaml}
records the upper edge $5.5$~TeV as the active budget, with historical edges
(CMS~2019 at $4.55$~TeV, PDG~Live~$3.8$~TeV, ATLAS~2013 at $2.07$~TeV,
ATLAS~2012 at $1.5$~TeV) retained for cross-check. \textbf{This is the one
constraint that fixes the smoke-test scale}: a model with $\MKK = 1$~TeV is below
the $5.5$~TeV edge and is \emph{excluded} by \code{CR001}; a model with $\MKK
\gtrsim 30$~TeV clears it (and every other \code{CR0xx} edge) easily.

\section{The KK graviton: diboson and diphoton (CR007, CR013)}
\label{sec:graviton}

The KK graviton is the genuinely RS-distinguishing state: it is spin 2, a direct
consequence of gravity propagating in the bulk. In the ``bulk RS'' scenario its
couplings are controlled by the dimensionless ratio $\tilde k = k/M_{\rm Pl}$,
and it decays preferentially to the longitudinal electroweak bosons and to the
Higgs/top --- i.e.\ to $WW$, $ZZ$ (\code{CR007}, final state \code{"WW, ZZ"}) and
to $\gamma\gamma$ (\code{CR013}, final state \code{"gamma gamma"}). The diphoton
channel is experimentally golden because it is narrow and clean, even though the
$\gamma\gamma$ branching fraction is small.

Our \code{CR007} active budget is the PDG-2025 bulk-graviton diboson edge at
$\tilde k = 0.5$ (\code{value\_id} ending \code{bulk\_graviton\_diboson\_mass\_kMpl\_0p5}),
$m(\GKK) > 1.4$~TeV; the YAML also carries stronger diphoton/dilepton edges (up
to $4.8$~TeV) as alternative benchmarks. \code{CR013} uses the CMS-2024 diphoton
RS-graviton benchmark at $\tilde k = 0.1$, $m_{G^\star} > 4.8$~TeV. The two
constraints read the same KK-graviton mass scale but compare it to
\emph{different} couplings ($\tilde k$) benchmarks --- a reminder that the
``mass edge'' is always conditional on an assumed coupling.

\section{KK electroweak vectors: diboson, dilepton, charged current (CR012, CR005, CR006)}
\label{sec:ewkk}

The KK partners of the electroweak gauge bosons are spin-1 resonances. A neutral
one decays to $WW/WZ/ZZ$ (\code{CR012}, treated through the heavy-vector-triplet
``HVT model~B'' $W'\to WZ$ benchmark, $m(W') > 4.4$~TeV) and to charged-lepton
pairs $\ell^+\ell^-$ (\code{CR005}, $Z'$-like dilepton, $\sim 5.15$~TeV); a
charged one is a $W'$-like state decaying to $\ell\nu$ or $tb$ (\code{CR006},
$\sim 6.0$~TeV). The dilepton and charged-current channels are the most powerful
\emph{if} the KK vector couples to leptons --- which in anarchic RS it does, but
the strength is model-dependent.

This is exactly why \code{CR005} and \code{CR006} are held \emph{out} of the
quark-only scan: they are \emph{lepton-involving} searches, and a quark-only
parameter point has no swept lepton-sector data to define the leptonic
couplings. Concretely, the \code{CR005} dilepton rate needs the KK-$Z'$ coupling
to electrons, $g_{ee}$, which is fixed by the lepton bulk-mass parameters
($c_L$, $c_e$) that a \emph{quark}-only scan never varies; with no such value in
the point's extras, recasting \code{CR005} would have to invent the coupling. So
the module is held out rather than fed a fabricated input. \code{CR012}
($WW/WZ/ZZ$, no final-state leptons required) reads only the geometry/quark
extras and stays \emph{in}. The split is bookkeeping about what data a point
carries, not a physics claim that leptonic channels are unimportant.

\section{Custodial vector-like quarks: pair production (CR002, CR003, CR004, CR008, CR010)}
\label{sec:vlq}

In the custodial RS model (\S\ref{part:model}) the third-generation quark
doublet is embedded in an enlarged representation, which brings along
\emph{vector-like quark} (VLQ) partners: new heavy quarks whose left- and
right-handed components transform identically (hence ``vector-like''), so their
mass is not tied to electroweak symmetry breaking. The custodial structure
predicts an \emph{exotic-charge} partner $T_{5/3}$ (electric charge $+5/3$), a
charge-$2/3$ partner $T$, a charge-$-1/3$ partner $B$, and weak-isospin
doublets. These are pair-produced by QCD and cascade to third-generation final
states:
\begin{center}
\begin{tabular}{lll}
\toprule
Constraint & Partner & Decay cascade \\
\midrule
\code{CR002} & $T_{5/3}$ (exotic charge $+5/3$) & $tW\,tW$ \\
\code{CR003} & $T$ (custodial, charge $2/3$) & $Wb/Zt/Ht$ \\
\code{CR004} & $B$ (charge $-1/3$) & $tW/bZ/bH$ \\
\code{CR008} & $T$ (singlet, non-custodial) & $Wb/Ht/Zt$ \\
\code{CR010} & $(T,B)$ weak-isospin doublet & mixed third-generation \\
\bottomrule
\end{tabular}
\end{center}
The exclusions are all $\sim 1.4$--$1.6$~TeV (e.g.\ \code{CR002} reads $T_{5/3}
> 1.46$~TeV from ATLAS; \code{CR003}'s active budget reads $m_T > 1.48$~TeV from
the CMS all-third-generation decay-mixture limit, CMS-B2G-20-011 /
arXiv:2209.07327; \code{CR004} reads $m_B > 1570$~GeV; \code{CR008} reads $>
1.36$~TeV; \code{CR010} reads $m_T = m_B > 1.37$~TeV). The stronger pure-$Wb$
ATLAS edge ($m_T > 1.70$~TeV, arXiv:2401.17165) is deliberately demoted to a
diagnostic in \code{CR003.py}: because a \code{ParameterPoint} carries no VLQ
branching fractions, the active HARD budget is the all-decay-mixture limit, not
the channel-pure one. These are markedly weaker than the gauge-resonance bounds because
pair production needs \emph{twice} the resonance mass in collision energy, so the
kinematic reach is roughly halved. \textbf{Physically, a VLQ partner mass is set
by the same KK scale, so in RS these searches probe $\MKK$ too} --- which is
exactly why our proxy interprets the scan's $\MKK$ as the VLQ mass
(\S\ref{sec:codeproxy}).

\section{The tails: contact and VBS and four-top (CR009, CR011, CR014)}
\label{sec:tails}

Two of the fourteen processes are genuinely \emph{not} on-shell resonance
searches. \code{CR009} is the high-mass Drell--Yan contact-interaction bound
(\S\ref{sec:eft}): it interprets the scan's KK scale as the EFT contact scale
$\Lambda_{\rm RS}$ and compares to the conservative (weakest) $\ell\ell qq$
compositeness limit ($\Lambda_{LL} > 23.9$~TeV, CMS destructive). \code{CR011}
is the longitudinal same-sign $WW$ VBS fiducial cross-section measurement (ATLAS
2025, $\sigma_{\rm fid} < 0.45$~fb at 95\% C.L.). A third lepton-involving
process, \code{CR014}, is grouped with these for allowlist purposes but is in
fact a \emph{resonance} search: BSM four-top production with a top-philic vector
mediator (CMS-B2G-25-005, $Z'\to t\bar t$ excluded up to $850$~GeV at
$\Gamma/m = 50\%$). It is held out not because it is a tail but because its
two-lepton selection makes it lepton-involving. All three are held \emph{out} of
the quark-only allowlist (\S\ref{sec:allowlist}).

\section{Where the near-flat \texorpdfstring{$\MKK$}{M\_KK} floor comes from}
\label{sec:floor}

Step back from the catalogue. Each search returns a mass edge $M_\star$ of order
a few TeV. The KK gluon gives the highest, $\sim 5.5$~TeV; the KK EW vectors and
gravitons give $\sim 4$--$6$~TeV (in their leptonic/diphoton channels) or
$\sim 1.4$~TeV (diboson); the VLQ partners give $\sim 1.4$--$1.7$~TeV. Taken as a
\emph{set}, and interpreting each edge as a lower bound on the common scale
$\MKK$, the union is a \emph{near-flat direct-search floor}: whatever the precise
benchmark, a model with $\MKK$ below a few TeV is excluded by \emph{something},
and a model with $\MKK$ well above $\sim 6$~TeV clears the quark-only IN
direct mass-edge set. The full catalogue still contains held-out tail/contact
entries such as \code{CR009}, so this sentence is not a statement about every
\code{CR0xx} module. This is the
``floor'' language of the orchestration ledger: at $\MKK = 1$~TeV \code{CR001}
(and most others) excludes; ``by $\sim 30$~TeV all clear.''

The floor is ``near-flat'' rather than a single sharp number precisely because it
is a \emph{union of unlike edges}, not one measurement repeated. The same scan
point $\MKK$ is judged against very different benchmarks at once: a resonance
edge (\code{CR001}, $5.5$~TeV), a pair-production edge (\code{CR002}, $1.46$~TeV
--- lower because pair production needs $\sim 2\MKK$ of collision energy), a
graviton edge (\code{CR007}/\code{CR013}, $1.4$/$4.8$~TeV, carrying the spin-2
root, \S\ref{sec:Mconvention}), and a contact-scale edge (\code{CR009},
$23.9$~TeV, a different observable entirely). The envelope of these dissimilar
edges is what looks ``flat'' from a distance; up close it is a staircase of
different physics. Below all of them the point is excluded by whichever bites
first; above the highest the whole set clears together. The reach caps where it
does for the parton-luminosity reason of the Orientation: at $\sqrt s = 13$~TeV
the cross section to make a multi-TeV resonance is suppressed by the steeply
falling parton densities, so no edge climbs much past a few TeV (the contact tail
excepted, since it bounds a virtual scale, not an on-shell mass).

\begin{callout}{Worked example --- the smoke test, one number at a time ($\MKK = 1$~TeV vs.\ $30$~TeV).}
Take a scan point with physical $\MKK = 1000$~GeV. The shared adapter converts
this to TeV, $m_{\rm pred} = 1.0$~TeV (\code{kk\_mass\_tev\_from\_m\_kk\_gev}
divides by $1000$). \code{CR001} loads the active edge $M_\star = 5.5$~TeV and
forms
\[
\text{ratio} = \frac{M_\star}{m_{\rm pred}} = \frac{5.5}{1.0} = 5.5 > 1
\quad\Rightarrow\quad \textbf{fails (vetoed)}.
\]
The pass test is literally \code{prediction.mass\_tev >= limit.value}, i.e.\
$1.0 \ge 5.5$ is False. Now take $\MKK = 30{,}000$~GeV, $m_{\rm pred} = 30$~TeV:
\[
\text{ratio} = \frac{5.5}{30} \approx 0.18 < 1 \quad\Rightarrow\quad \textbf{passes}.
\]
Every other \code{CR0xx} edge ($23.9$~TeV for the conservative contact bound,
$\le 6$~TeV for the resonances) is likewise cleared at $30$~TeV. This single pair of numbers
is the whole behaviour: the direct floor turns off somewhere between a few TeV
and (for the contact tail) tens of TeV, and the rigorous flavour floor sits
\emph{above} that --- which is why these constraints are present for completeness
but rarely decisive.

\medskip
\emph{The same $\MKK$, unlike edges.} The $1$~TeV smoke test above hides the more
confusing point that one scan value is judged against \emph{dissimilar} edges at
once. Take the borderline $m_{\rm pred} = 2.0$~TeV and run three processes:
\[
\underbrace{\tfrac{5.5}{2.0}=2.75>1}_{\text{\code{CR001} resonance: FAILS}},
\qquad
\underbrace{\tfrac{1.46}{2.0}=0.73<1}_{\text{\code{CR002} VLQ pair: PASSES}},
\qquad
\underbrace{\tfrac{23.9}{2.0}=11.95>1}_{\text{\code{CR009} contact: FAILS}}.
\]
The \emph{same} $\MKK$ clears the pair-production edge (which is low because pair
production costs $\sim 2\MKK$ of collision energy) while failing both the
resonance edge and the contact-scale edge --- this is the ``union of unlike
edges'' of \S\ref{sec:floor} made arithmetic, and it is why the floor is a
staircase, not one number.

\medskip
\emph{The $2.45$ (and $3.83$) convention in action.} Suppose the scan instead
reports $\LIR = 2.0$~TeV (i.e.\ the $\MKK\equiv\LIR$ convention) rather than a
physical mass. To compare against a collider edge --- always quoted in
\emph{physical} resonance mass --- one must first apply the Bessel root: for the
KK-gluon resonance (\code{CR001}) the physical mass is
$x_1^{\rm gauge}\LIR \approx 2.45\times 2.0 = 4.9$~TeV, whereas for the spin-2
graviton (\code{CR007}/\code{CR013}) it is $x_1^{\rm grav}\LIR \approx 3.83\times
2.0 = 7.66$~TeV. Feeding the bare $2.0$~TeV into either comparison --- or feeding
$2.45$ to the graviton --- mis-scales the ratio. This is the conversion
\S\ref{sec:Mconvention} calls mandatory; the proxy is only consistent if the
scan and the edge already share one convention.
\end{callout}

% ====================================================================
\part{Model dependence, and the honest ``proxy'' label}
\label{part:model}
% ====================================================================

\section{Custodial symmetry: which searches even exist}
\label{sec:custodial}

Custodial symmetry --- the enlarged bulk gauge group $\SUL\times\SUR\times\UX$
that tames the oblique $T$ parameter --- changes the collider menu in a concrete
way: it \emph{introduces} the vector-like quark partners. The exotic-charge
$T_{5/3}$ (\code{CR002}) is a smoking-gun of the custodial bidoublet embedding of
the top; it simply does not exist in the minimal model. So \code{CR002}--\code{CR004}
and \code{CR010} are \emph{custodial-specific} searches, while \code{CR008} (a
non-custodial singlet $T$) is their minimal-model counterpart. The sidecars label
this explicitly (\code{subclass: "Custodial fermion partner; charge 2/3"} for
\code{CR003}).

The flip side: custodial symmetry \emph{lowers} the indirect floor (it removes
the volume-enhanced $T$ and the $Zb_L$ veto), which is precisely the regime where
a few-TeV direct floor could in principle matter. In the custodial scan the
rigorous floor drops to the $\Delta F = 2$ set ($\sim 2$--$3$~TeV), and the
inclusive floor ($\sim 7$~TeV) is then set by the \emph{proxy} oblique-$S$ plus
these \code{CR*} searches. So the direct collider bounds are most relevant
exactly when custodial protection is on --- but even then they sit at the same
order as, and are encoded less rigorously than, the proxy oblique-$S$.

\section{Flavour alignment: does it move the collider bound?}
\label{sec:alignment}

Flavour alignment (aligning the bulk masses to suppress flavour-changing neutral
currents) is the standard way to \emph{lower} the indirect $\Delta F = 2$ floor.
It does \emph{not} directly relax a direct collider bound, because the production
of a KK gluon or graviton is governed by its (flavour-diagonal) coupling to light
quarks and gluons, not by the flavour-violating couplings that alignment
suppresses. What alignment \emph{can} change is the decay pattern (e.g.\ the
top-philic branching fractions) and hence the precise mass edge --- but our proxy
does not use branching fractions at all, so in our encoding alignment has
\emph{no effect} on the direct floor. This is a feature of the proxy's
crudeness, not a physics statement, and it is worth flagging as such.

\section{Is our treatment rigorous? No --- it is a mass-exclusion proxy}
\label{sec:proxy}

Here is the blunt assessment, and it is the same for all fourteen processes. A
\emph{rigorous} recast of any of these searches would, for each scan point,
compute the production cross section $\sigma(pp\to X)$ from the RS couplings, the
branching fractions into the searched channel, the total width and line shape,
the interference with the SM continuum, fold all of that through the published
analysis acceptance, and compare to the experiment's \emph{mass-dependent} limit
curve. Our code does \emph{none} of this. It takes a single KK mass scale from
the scan point, calls it the resonance (or VLQ, or contact-scale) mass, and asks
only whether that mass clears the catalogued benchmark edge. Every module carries
the assumption string \code{COLLIDER\_RESONANCE\_MASS\_PROXY\_ASSUMPTION\_V1}
(or a process-specific analogue) whose text begins
\code{"NEEDS-HUMAN-PHYSICS: ... uses the catalogued benchmark mass-exclusion edge
as a mass lower bound. It does not compute sigma(pp->X)*BR(X->final state),
width, interference, or acceptance."} The matching status in the module
docstrings is uniformly \code{NEEDS-HUMAN-PHYSICS}; the relevance class in the
sidecars is \code{"Cross-check ... subleading to low-energy flavor constraints"}.
\textbf{Tag honestly: every \code{CR0xx} is a \emph{proxy}, not rigorous.}

% ====================================================================
\part{What is in OUR code}
\label{part:code}
% ====================================================================

\emph{Reader here for the physics may skim Part~\ref{part:code} on a first pass
and jump to the subtleties in Part~\ref{part:subtle}; this part documents the
exact code so the note doubles as an implementation reference.}

\section{The shared comparison machinery}
\label{sec:codeproxy}

The reusable core lives in \code{quarkConstraints/collider\_resonance.py}, reached
only through the boundary adapter
\code{flavor\_catalog\_constraints/physics\_adapters/collider\_resonance.py}.
Three dataclasses carry the logic:
\code{ColliderResonancePrediction} (the model side: \code{resonance},
\code{final\_state}, \code{mass\_tev}, optional \code{sigma\_times\_br}),
\code{ColliderResonanceLimit} (the experiment side: \code{value}, \code{units},
\code{limit\_kind}, \code{benchmark\_model}, $\dots$), and
\code{ColliderResonanceComparison} (the result). The mass conversion is the
one-liner
\begin{equation}
\code{kk\_mass\_tev\_from\_m\_kk\_gev}:\quad m_{\rm TeV} = m_{\rm GeV}/1000,
\label{eq:gev2tev}
\end{equation}
guarded to be positive and finite. The comparison is
\code{evaluate\_resonance\_limit}, whose two branches are
\begin{align}
\text{(mass lower bound)}\quad
&\text{ratio} = \frac{m_{\rm limit}}{m_{\rm pred}}, \quad
\text{passes} \iff m_{\rm pred} \ge m_{\rm limit};
\label{eq:ratiomass}\\
\text{($\sBR$ upper limit)}\quad
&\text{ratio} = \frac{\sBR_{\rm pred}}{\text{limit}}, \quad
\text{passes} \iff \text{ratio} \le 1.
\label{eq:ratiosbr}
\end{align}
The two constants \code{MASS\_LOWER\_BOUND = "mass\_lower\_bound"} and
\code{SIGMA\_TIMES\_BR\_UPPER\_LIMIT = "sigma\_times\_br\_upper\_limit"} name the
two branches. In production \emph{every} \code{CR0xx} uses the first branch
(mass): the second exists for a future $\sBR$ recast but is never reached,
because no scan point carries a \code{sigma\_times\_br} prediction. The
comparison also enforces that the prediction's \code{resonance} and
\code{final\_state} strings match the limit's --- a guard against, say, comparing
a $t\bar t$ prediction to a $\gamma\gamma$ limit.

\section{Per-process mass resolution and the extras they read}
\label{sec:extras}

Each \code{CR0xx.py} resolves a mass from the \code{ParameterPoint} ``extras''
and builds a prediction through a process-specific helper in the adapter. The
resolution order is documented and deterministic. For \code{CR001},
\code{resolve\_kk\_gluon\_mass\_gev} reads the explicit \code{kk\_gluon\_mass\_gev}
extra if present, else falls back to \code{quark\_mass\_basis\_couplings.M\_KK}
(the function \code{mass\_from\_source\_gev} pulls the \code{M\_KK} attribute).
The electroweak-sector processes read \code{kk\_ew\_mass\_gev} first, falling
back to \code{M\_KK}. The proxy-assumption strings make the interpretation
explicit, e.g.\ for the VLQ searches \code{VLQ\_PAIR\_MASS\_PROXY\_ASSUMPTION\_V1}
declares \code{"uses m\_VLQ = M\_KK as a KK-fermion mass proxy"}, and the
charge-$2/3$/$-1/3$ variants (\code{VLQ\_T\_PAIR\_...}, \code{VLQ\_B\_PAIR\_...})
analogously set $m_T = \MKK$, $m_B = \MKK$. The contact-scale resolver
\code{resolve\_dy\_contact\_scale\_gev} (CR009) prefers \code{kk\_ew\_mass\_gev},
then \code{kk\_gluon\_mass\_gev}, then \code{M\_KK}, as the
$\Lambda_{\rm RS}$ proxy. The full map of which extras each constraint reads is:
\begin{center}\small
\begin{tabular}{lll}
\toprule
Constraint & Mass proxy interpretation & Extras read \\
\midrule
\code{CR001} & $m(\gKK) = $ \code{kk\_gluon\_mass\_gev} or \code{M\_KK} & \code{kk\_gluon\_mass\_gev}, couplings \\
\code{CR002/3/4} & $m_{\rm VLQ} = \MKK$ & \code{quark\_mass\_basis\_couplings} \\
\code{CR007} & $m(\GKK) = $ \code{kk\_ew}/\code{kk\_gluon}/\code{M\_KK} & all three \\
\code{CR008/10} & $m_{\rm VLQ} = \MKK$ & \code{quark\_mass\_basis\_couplings} \\
\code{CR012} & $m_{V_{\rm KK}} = $ \code{kk\_ew\_mass\_gev} or \code{M\_KK} & \code{kk\_ew\_mass\_gev}, couplings \\
\code{CR013} & $m(\GKK) = $ \code{kk\_ew}/\code{kk\_gluon}/\code{M\_KK} & all three \\
\midrule
\multicolumn{3}{l}{\emph{held OUT of quark-only (lepton-involving / INFO):}} \\
\code{CR005} & $m(Z'_{\rm KK}) = $ \code{kk\_ew} or \code{M\_KK} & \code{kk\_ew\_mass\_gev}, couplings \\
\code{CR006} & $m(W'_{\rm KK}) = $ \code{kk\_ew} or \code{M\_KK} & \code{kk\_ew\_mass\_gev}, couplings \\
\code{CR009} & $\Lambda_{\rm RS} = $ \code{kk\_ew}/\code{kk\_gluon}/\code{M\_KK} & all three \\
\code{CR011} & (VBS fiducial $\sigma$; no mass) & none \\
\code{CR014} & $m(Z'_{\rm top\text{-}philic}) = $ \code{kk\_ew} or \code{M\_KK} & \code{kk\_ew\_mass\_gev}, couplings \\
\bottomrule
\end{tabular}
\end{center}
These extras are read from \code{QUARK\_ONLY\_ALLOWLIST\_EXTRAS} in the scan
machinery; the quark-only point already supplies \code{kk\_gluon\_mass\_gev},
\code{kk\_ew\_mass\_gev}, and \code{quark\_mass\_basis\_couplings}, so no silent
SM/zero default is injected.

\section{Severities and how a point is vetoed}
\label{sec:severity}

Thirteen of the fourteen constraints are \code{Severity.HARD}: a point whose
predicted mass falls \emph{below} the active benchmark edge fails and is
\emph{vetoed}. The reported \code{ratio} is $m_{\rm limit}/m_{\rm pred}$, so
$\text{ratio} \le 1$ passes, matching the catalogue contract that ``ratio above
budget'' is a failure. The lone exception is \code{CR011} (VBS), which is
\code{Severity.INFO}: it forms the advisory comparison $\sigma_{\rm fid}/0.45\,
\text{fb} \le 1$ but is \emph{non-vetoing} by construction --- its own code notes
\code{"INFO/non-vetoing even when advisory passes=False"}, and because no scan
point supplies a predicted fiducial $W_L W_L$ cross section, it records a
non-vetoing result without ever recasting. So \code{CR011} is informational only:
it never excludes a point.

\section{Missing-input behaviour (fail-safe vs.\ fail-loud)}
\label{sec:missing}

The modules distinguish two failure modes carefully. If the mass extras are
entirely \emph{absent}, the constraint returns \code{passes=True} with a note
that it ``was not evaluated'' (it does not veto on missing data --- see
\code{CR001.evaluate} returning early when \code{m\_kk\_gev is None}). If the
extras are \emph{present but invalid} (non-numeric, non-finite, non-positive),
the constraint returns \code{passes=False} with an \code{exception\_type}
diagnostic --- it fails loudly rather than silently passing a malformed point.
This asymmetry is deliberate: absence of data is not evidence of exclusion, but
corrupt data should not slip through.

\section{The active benchmark, exactly}
\label{sec:active}

Each YAML stores a \code{pdg\_or\_equivalent.values} list of historical and
current edges; the Python pins one as the \emph{active} budget by its
\code{value\_id}. For \code{CR001} the active id is
\code{CMS2026:CR001:gkk\_ttbar\_mass\_exclusion} ($5.5$~TeV, CMS-B2G-25-009);
the historical CMS-2019/PDG/ATLAS edges are loaded too but only reported in
\code{diagnostics["historical\_mass\_limits\_tev"]}. The budget policy string is
explicit: \code{"active HARD budget is the CMS2026 current canonical RS KK-gluon
ttbar benchmark mass-exclusion upper edge ... interpreted as a lower bound above
the excluded interval for that benchmark."} A representative active-edge map:
\begin{center}\small
\begin{tabular}{llll}
\toprule
Constraint & Active edge & Experiment / source & Final state \\
\midrule
\code{CR001} & $5.5$~TeV & CMS-B2G-25-009, arXiv:2603.23454 & $t\bar t$ \\
\code{CR002} & $1.46$~TeV & ATLAS, PDG q009 / arXiv:2212.05263 & $tW\,tW$ \\
\code{CR003} & $1.48$~TeV & CMS-B2G-20-011, arXiv:2209.07327 & $Wb/Zt/Ht$ \\
\code{CR004} & $1.57$~TeV & CMS, PRD~110 (2024) 052004 & $tW/bZ/bH$ \\
\code{CR007} & $1.4$~TeV ($\tilde k{=}0.5$) & PDG-2025 bulk graviton & $WW/ZZ$ \\
\code{CR008} & $1.36$~TeV & ATLAS, PDG t$'$ AAD~24N & $Wb/Ht/Zt$ \\
\code{CR010} & $1.37$~TeV & ATLAS Run-2 VLQ comb. & mixed 3rd-gen \\
\code{CR012} & $4.4$~TeV & CMS-B2G-20-009 (HVT~B $W'{\to}WZ$) & $WW/WZ/ZZ$ \\
\code{CR013} & $4.8$~TeV ($\tilde k{=}0.1$) & CMS-2024 RS-graviton diphoton & $\gamma\gamma$ \\
\midrule
\multicolumn{4}{l}{\emph{held OUT of quark-only:}} \\
\code{CR005} & $5.15$~TeV & CMS, JHEP~07 (2021) 208 & $\ell^+\ell^-$ \\
\code{CR006} & $6.0$~TeV & ATLAS, arXiv:1906.05609 & $\ell\nu$ \\
\code{CR009} & $23.9$~TeV ($\Lambda_{LL}$, conservative) & CMS, arXiv:2103.02708 & $\ell\ell qq$ \\
\code{CR011} & $0.45$~fb (INFO) & ATLAS, PRL~135 (2025) 111802 & $W_L W_L$ VBS \\
\code{CR014} & $0.85$~TeV & CMS-B2G-25-005 ($\Gamma/m{=}50\%$) & four-top \\
\bottomrule
\end{tabular}
\end{center}
(Values are the active budgets as stored in the sidecars; the YAMLs also carry
alternative/historical edges for auditability.)

\section{The quark-only allowlist split: 9 IN, 5 OUT}
\label{sec:allowlist}

The change recorded in \code{docs/quark\_scan\_constraint\_update\_2026-06.md}
(Change~2) added the \emph{lepton-free} collider searches to the quark-only
allowlist, growing it from $37$ to $46$ constraints. The split is encoded in the
scan machinery (\code{QUARK\_ONLY\_ALLOWLIST\_EXTRAS} and the IN/OUT id lists):
\begin{itemize}[leftmargin=1.4em]
\item \textbf{IN (9, proxy / HARD):} \code{CR001} ($\gKK\to t\bar t$);
\code{CR002}, \code{CR003}, \code{CR004}, \code{CR008}, \code{CR010} (VLQ pair
production); \code{CR007} \& \code{CR013} (KK graviton $\to WW/ZZ$,
$\gamma\gamma$); \code{CR012} (KK spin-1 $\to WW/WZ/ZZ$). Each was verified to
read only quark/geometry extras (\code{kk\_gluon\_mass\_gev},
\code{kk\_ew\_mass\_gev}, \code{quark\_mass\_basis\_couplings}).
\item \textbf{OUT (5, deferred, lepton-involving or INFO):} \code{CR005}
(dilepton resonance), \code{CR006} ($W'\to\ell\nu$), \code{CR009} ($\ell\ell qq$
contact), \code{CR011} (VBS, INFO/non-vetoing), \code{CR014} (four-top with
two-lepton selection).
\end{itemize}
The logic is purely about \emph{what data a quark-only point carries}: a point
with no swept lepton sector cannot define the leptonic couplings that the OUT
searches need, so including them would invent inputs. \code{CR011} is also OUT
because it is INFO and would never veto anyway. The net effect, per the ledger:
at $\MKK = 1$~TeV the excluding constraints include the rigorous $\Delta F = 2$
set ($\eps_K = $ \code{K001}, $\Delta m_s = $ \code{B003}, $\phi_s = $
\code{B004}) and the now-live $Z\to b\bar b$ (\code{T010}/\code{T011}, rigorous)
\emph{plus} the collider direct searches (\code{CR001} proxy, $\dots$); by
$\sim 30$~TeV all clear.

% ====================================================================
\part{Subtleties and the gap to ``rigorous''}
\label{part:subtle}
% ====================================================================

\section{The proxy compares masses, not rates}
The single defining limitation: \eqref{eq:ratiomass} compares the scan's KK mass
to an experimental \emph{mass edge} that was itself derived under a specific
benchmark signal model. If the real RS point has a smaller production cross
section than the benchmark (e.g.\ heavier/more-UV-localised partons feeding the
resonance), a \emph{lighter} resonance could survive that the proxy would
exclude --- and vice versa. The proxy is therefore neither conservative nor
aggressive in a controlled way; it is simply a mass cut at a benchmark-defined
edge. This is the content of the uniform \code{NEEDS-HUMAN-PHYSICS} tag.

\section{The \texorpdfstring{$\MKK$}{M\_KK} convention: the factor of $2.45$ again}
\label{sec:Mconvention}
``$\MKK$'' is overloaded in exactly the way the $S,T,U$ note flags for the
oblique sector. The collider papers report the mass of a \emph{physical}
resonance --- $\MKK^{\rm phys} = x_1\LIR$ \eqref{eq:KKmass}. But the scan's
reported $\MKK$, and any theory formula written with $\MKK\equiv\LIR$, differ by
that factor $x_1$. Crucially, \emph{$x_1$ is not universal across the
processes}: for the KK gauge and VLQ states it is $x_1^{\rm gauge}\approx 2.45$,
whereas for the spin-2 KK \emph{gravitons} (\code{CR007}, \code{CR013}) it is
$x_1^{\rm grav}\approx 3.83$ --- so a single physical $\LIR$ feeds a graviton
edge that is $\sim 1.5\times$ higher in physical mass than the gauge edge. The
ledger spells out the gauge-sector consequence for the \emph{rigorous} $Z\to
b\bar b$ floor: ``$\sim 25$--$30$~TeV physical $\leftrightarrow \sim 10$--$12$~TeV
in $\LIR$ units'' (this $\sim 2.45$ mapping is correct for gauge/$Zb\bar b$, but
the graviton constraints would carry $\sim 3.83$). \textbf{The same care is
mandatory here}: whatever scale the scan feeds into \eqref{eq:gev2tev} must be in
the \emph{same} convention as the experimental edge it is compared to, and for
the graviton processes the \emph{spin-2 root} must be used, not the gauge one. If
the scan passes a physical $\MKK$ but the benchmark edge is quoted (as collider
edges are) in physical resonance mass, the comparison is consistent; mixing
conventions --- or applying the gauge root to a graviton --- would mis-scale the
\code{CR0xx} ratio (by $\sim 2.45$ for gauge, $\sim 3.83$ for gravitons).

\section{Pair production halves the reach}
The VLQ searches (\code{CR002}--\code{CR004}, \code{CR008}, \code{CR010}) are
pair-production bounds, so their mass edges ($\sim 1.4$--$1.7$~TeV) are
\emph{not} directly comparable to the single-resonance edges ($\sim 4$--$6$~TeV):
pair production needs $\sim 2 M_Q$ in partonic energy. Interpreting both as
``lower bounds on the same $\MKK$'' (as the proxy does) is a crude
over-simplification --- the VLQ partner mass and the KK gauge mass are both set
by $\LIR$ but with different $\mathcal{O}(1)$ coefficients. A rigorous treatment
would carry separate spectra; the proxy collapses them to one $\MKK$.

\section{Benchmark coupling dependence ($\tilde k$, $\Gamma/M$)}
Several edges are explicitly conditional on a coupling benchmark: \code{CR007}
uses $\tilde k = k/M_{\rm Pl} = 0.5$, \code{CR013} uses $\tilde k = 0.1$,
\code{CR014} uses $\Gamma/m = 50\%$, and \code{CR001}'s line shape assumes
$\Gamma/M \approx 0.153$. The active edge moves with these choices (the YAMLs
carry alternative-benchmark values). Our proxy pins one benchmark per process and
does not propagate the model point's \emph{actual} coupling into the edge --- a
deferred piece.

\section{Tail observables shoe-horned into a mass cut}
\code{CR009} (a contact scale $\Lambda$) and \code{CR011} (a fiducial VBS cross
section) are not resonance masses, yet \code{CR009} is forced through the same
mass-lower-bound comparison by treating $\Lambda_{\rm RS}$ as a ``mass.'' That is
dimensionally fine (both are TeV scales) but physically a different observable;
the relation $\Lambda \sim \MKK/\sqrt{|g_q g_\ell|}$ \eqref{eq:contact} carries
coupling factors the proxy ignores. \code{CR011} side-steps the issue by being
INFO/non-vetoing.

\section{The rigorous floor smothers the direct floor}
The practical reason these constraints rarely decide anything: in the anarchic
quark scan the rigorous flavour bounds already push the physical $\MKK$ floor to
$\sim 25$--$30$~TeV (\code{T010}/\code{T011} $Z\to b\bar b$ plus $\Delta F = 2$),
and the methodology note records even higher anarchic comparison scales
(\code{CR001.yaml} auxiliary inputs: RUNA p50 $\MKK^{\rm min} = 47.26$~TeV,
p95 $127.13$~TeV at $g_{s\star} = 3$). A few-TeV direct floor cannot bite above
that. The direct searches matter only in the \emph{custodial} regime, where the
rigorous floor drops to $\sim 2$--$3$~TeV; even there the proxy oblique-$S$ and
the \code{CR*} proxies set the inclusive ($\sim 7$~TeV) floor, and both are
\emph{proxy}, not rigorous --- so the honest headline custodial floor remains the
rigorous $\Delta F = 2$ value.

\section{Checklist: what ``rigorous direct bounds'' would require}
\begin{enumerate}[leftmargin=1.8em]
\item Replace the mass-edge comparison \eqref{eq:ratiomass} with a per-point
$\sBR$ computation (the coded-but-unused branch \eqref{eq:ratiosbr}), using RS
production couplings and branching fractions.
\item Fold in width, line shape, interference, and the published analysis
acceptance; compare to the experiment's \emph{mass-dependent} limit curve, not a
single edge.
\item Carry separate KK-gauge and VLQ-partner spectra rather than one $\MKK$, and
fix a single documented $\MKK$ convention (physical vs.\ $\LIR$).
\item Propagate the model point's actual coupling benchmark ($\tilde k$,
$\Gamma/M$) into the chosen edge.
\item Treat the contact (\code{CR009}) and VBS (\code{CR011}) tails as the EFT /
fiducial observables they are, not as mass cuts.
\end{enumerate}

\begin{callout}{Common confusions, defused.}
\begin{itemize}[leftmargin=1.2em,nosep]
\item[\textbf{A.}] \emph{``The direct searches are our strongest bound because
they probe real particles.''} No --- they are conceptually the cleanest but
\emph{numerically the weakest} here. The 13~TeV collider caps the reach at a few
TeV; the indirect flavour bounds already sit at $\sim 25$--$30$~TeV (physical).
\item[\textbf{B.}] \emph{``Our code computes the cross section.''} No --- it
compares \emph{masses} to a benchmark edge. No $\sigma$, BR, width,
interference, or acceptance is computed; every module is tagged
\code{NEEDS-HUMAN-PHYSICS}.
\item[\textbf{C.}] \emph{``$\MKK$ means one thing.''} Overloaded by the Bessel
root $x_1$ (physical first KK mass vs.\ $\LIR$, \S\ref{sec:Mconvention}), and
$x_1$ is \emph{spin-dependent}: $x_1^{\rm gauge}\approx 2.45$ for KK gauge/VLQ
states but $x_1^{\rm grav}\approx 3.83$ for the spin-2 KK gravitons (\code{CR007},
\code{CR013}), so the graviton edges carry an extra factor $\sim 1.5$.
\item[\textbf{C$'$.}] \emph{``All the mass edges are the same kind of bound.''}
No --- resonance ($\sim 5$~TeV) and pair-production ($\sim 1.5$~TeV) edges differ
because pair production needs $\sim 2 M_Q$ in collision energy; the proxy
collapses both onto one $\MKK$ (\S\ref{sec:Mconvention}).
\item[\textbf{D.}] \emph{``CR011 (VBS) can veto a point.''} No --- it is
\code{Severity.INFO}, non-vetoing even when its advisory comparison returns
\code{passes=False}.
\item[\textbf{E.}] \emph{``All 14 run in the quark-only scan.''} No --- 9
lepton-free ones are IN, 5 lepton-involving/INFO (\code{CR005}, \code{CR006},
\code{CR009}, \code{CR011}, \code{CR014}) are OUT.
\end{itemize}
\end{callout}

% ====================================================================
\part*{Annotated reading guide}
\addcontentsline{toc}{section}{Annotated reading guide}
% ====================================================================

\noindent\emph{How to read this note.} A \textbf{physics-first} reader who knows
nothing should take Parts~\ref{part:phys}--\ref{part:rs} (the collider physics
and the RS spectrum) and then the ``common confusions'' box, skimming the
code-detail Part~\ref{part:code}. An \textbf{implementer} can jump straight to
Part~\ref{part:code} (modules, severities, active edges, allowlist) and use
Parts~\ref{part:phys}--\ref{part:rs} as background. Everyone should read
Part~\ref{part:subtle} before quoting a bound externally. The references below
are grouped by process.

\begin{description}[leftmargin=0em,style=nextline]
\item[Lillie, Randall, Wang, arXiv:hep-ph/0701166.] The theory baseline for the
KK gluon as a top-philic LHC resonance; motivates boosted-top reconstruction. The
paper-era reference in \code{CR001.yaml}; read alongside \S\ref{sec:cr001}. Not a
measured limit (a discovery-reach projection).

\item[CMS-B2G-25-009, arXiv:2603.23454.] The current best full-Run-2 CMS
all-topology $t\bar t$ resonance search; excludes the RS KK-gluon benchmark up to
$5.5$~TeV. The \emph{active} budget of \code{CR001} and the constraint that sets
the smoke-test scale.

\item[ATLAS $T_{5/3}$ search, arXiv:2212.05263; CMS-B2G-20-011, arXiv:2209.07327
(active \code{CR003}); ATLAS $T\to Wb$, arXiv:2401.17165 (\code{CR003}
diagnostic).] Representative vector-like-quark pair-production limits
($\sim 1.4$--$1.7$~TeV) feeding \code{CR002}/\code{CR003}. The active
\code{CR003} budget is the CMS all-decay-mixture edge ($1.48$~TeV); the stronger
ATLAS pure-$Wb$ edge ($1.70$~TeV) is held as a diagnostic. Read with
\S\ref{sec:vlq}; the exotic-charge $T_{5/3}$ is the custodial smoking gun.

\item[CMS high-mass diphoton RS graviton (CMS-2024); PDG-2025 Extra Dimensions
listing.] Source of the \code{CR013}/\code{CR007} KK-graviton edges
($\sim 4.8$~TeV diphoton at $\tilde k = 0.1$; $1.4$~TeV diboson at $\tilde k =
0.5$). The spin-2 state is the RS-distinguishing signature (\S\ref{sec:graviton}).

\item[CMS-B2G-20-009 (HVT $W'\to WZ$).] The heavy-vector-triplet diboson
benchmark behind \code{CR012} ($m(W') > 4.4$~TeV). The model-independent spin-1
diboson route (\S\ref{sec:ewkk}).

\item[PDG-2025 quark--lepton compositeness review (CMS, arXiv:2103.02708).] The
$\ell\ell qq$ contact-interaction bound behind \code{CR009}; the active budget is
the \emph{conservative} (weakest) destructive edge $\Lambda_{LL} > 23.9$~TeV, with
stronger benchmarks (ATLAS $\Lambda_{LL}^{+} > 35.8$~TeV) retained as diagnostics.
The EFT-tail observable of \S\ref{sec:eft}. Held OUT of quark-only.

\item[ATLAS longitudinal same-sign $WW$ VBS, PRL~135 (2025) 111802
(arXiv:2503.11317).] The fiducial cross-section measurement behind \code{CR011}
($< 0.45$~fb), our one INFO/non-vetoing collider record (\S\ref{sec:severity}).

\item[Repo code.] \code{quarkConstraints/collider\_resonance.py} (the shared
$\sBR$/mass machinery), the boundary adapter
\code{flavor\_catalog\_constraints/physics\_adapters/collider\_resonance.py}, the
fourteen \code{.../primary/collider\_rs/CR0xx.py}, their sidecars
\code{flavor\_catalog/processes/collider\_rs/CR0xx.yaml}, and the allowlist split
in \code{docs/quark\_scan\_constraint\_update\_2026-06.md} (Change~2) with
\code{QUARK\_ONLY\_ALLOWLIST\_EXTRAS}. \textbf{These are the source of truth for
every number and tag in this note.}
\end{description}

\bigskip
\noindent\rule{\linewidth}{0.4pt}\\
\footnotesize
\emph{Internal note. All active benchmark edges, severities, extras, and
``proxy vs rigorous'' tags are as stored in the \code{collider\_rs} sidecars and
modules and in the orchestration ledgers; every \code{CR0xx} is a
mass-exclusion \emph{proxy} (\code{NEEDS-HUMAN-PHYSICS}), not a rigorous recast.
The $\MKK$ convention (factor $x_1$, spin-dependent: $\approx 2.45$ for KK
gauge/VLQ, $\approx 3.83$ for KK gravitons) must be checked before quoting a
TeV bound externally.}

\end{document}
